\clearpage\section{研究方法}

\subsection{研究架構}

\subsubsection{共用資料層與三個實驗分支}

系統原有的協調者（Coordinator）、資料快照（Data Snapshot）與四類研究代理人（Research Agents）都保留。共用資料層會產生一份「不帶投資方向」的鎖定研究資料與證據登錄。RQ1、RQ2 讀的是同一份中立報告（Neutral Report）；RQ3 的通才組也讀這份整合報告，而領域組則使用同一份研究輸出依領域切開的部分，不另外抓新資料，也不用別的摘要模型重做輸入。整體概念架構如圖~\ref{fig:flow}。
\begin{figure}[H]
\centering
\begin{tikzpicture}[font=\footnotesize,>={Stealth[length=2mm]},
  box/.style={draw=black!55,rounded corners=3pt,align=center,inner sep=4pt,minimum height=8mm,fill=white},
  shared/.style={box,fill=blue!6,draw=blue!50!black,text width=13.4cm},
  done/.style={box,draw=teal!60!black,fill=teal!12,thick,text width=4.3cm,minimum height=21mm},
  plan/.style={box,draw=orange!75!black,fill=orange!8,thick,dashed,text width=4.3cm,minimum height=21mm},
  post/.style={box,draw=black!60,fill=black!4,text width=13.4cm},
  arr/.style={->,thick,black!65}]
% 共用資料層底框
\draw[rounded corners=5pt,draw=blue!40,fill=blue!3,dashed] (0.1,0.8) rectangle (14.9,-2.6);
\node[anchor=south west,font=\scriptsize\sffamily,text=blue!50!black] at (0.1,0.85) {共用資料層（三個分支讀同一份）};
\node[shared] (data) at (7.5,0) {資料來源：行情、財報、新聞標題、總經\\依申報日、歷史版本與時間戳記篩選};
\node[shared] (snap) at (7.5,-1.7) {不可變資料快照 $\rightarrow$ 四個研究代理人 $\rightarrow$ 中立報告與證據登錄};
\draw[arr] (data) -- (snap);
% 三個分支
\node[done] (rq1) at (2.5,-4.9) {\textbf{RQ1}（已完成）\\同一模型 M1/M1\\\grp{C} 對 \grp{D}\\讀同一份中立報告};
\node[plan] (rq2) at (7.5,-4.9) {\textbf{RQ2}（設計階段）\\不同模型 M1/M2\\\grp{C} 對 \grp{D}\\配置 X、Y 平衡};
\node[plan] (rq3) at (12.5,-4.9) {\textbf{RQ3}（設計階段）\\通才組 對 領域專業化組\\四領域各自多空辯論\\（規則 \grp{C}）};
\foreach \n in {rq1,rq2,rq3}{\draw[arr] (\n.north |- snap.south) -- (\n.north);}
% 裁決層
\node[box,text width=9.3cm,minimum height=11mm] (judge) at (5,-7.3) {裁決者（固定；逐項核對證據後擇一方向）};
\node[box,text width=4.3cm,minimum height=11mm] (agg) at (12.5,-7.3) {最終整合者\\（RQ3；本身即最上層裁決）};
\draw[arr] (rq1.south) -- (rq1.south |- judge.north);
\draw[arr] (rq2.south) -- (rq2.south |- judge.north);
\draw[arr] (rq3.south) -- (agg.north);
% 守門員與回測
\node[post] (gate) at (7.5,-9.0) {守門員：檢查證據引用、資料覆蓋、波動度與風險；必要時降級};
\draw[arr] (judge.south) -- (judge.south |- gate.north);
\draw[arr] (agg.south) -- (agg.south |- gate.north);
\node[post] (back) at (7.5,-10.5) {回測 $\rightarrow$ 準確率／Sharpe／成本／診斷};
\draw[arr] (gate) -- (back);
% 圖例
\node[draw=teal!60!black,fill=teal!12,thick,rounded corners=2pt,minimum width=9mm,minimum height=4mm] (l1) at (3.8,-11.6) {};
\node[anchor=west,font=\scriptsize] at (l1.east) {實線：已完成};
\node[draw=orange!75!black,fill=orange!8,thick,dashed,rounded corners=2pt,minimum width=9mm,minimum height=4mm] (l2) at (8.3,-11.6) {};
\node[anchor=west,font=\scriptsize] at (l2.east) {虛線：設計階段，尚未執行};
\end{tikzpicture}
\caption{研究概念架構。三個分支共用同樣的資料、候選決策與評估方式；實線為已完成，虛線為設計階段（尚未執行）。}
\label{fig:flow}
\end{figure}

\subsubsection{裁決者與守門員}

裁決者（Judge）綜合辯論內容、核對證據，並決定投資方向。最終整合者（Final Aggregator）只用在 RQ3，負責整合四個領域的看法；它本身就是最上層的裁決者，因此不再另外設一位決策裁決者。守門員（Gatekeeper）在投資方向決定之後才出場，檢查證據引用、資料覆蓋、波動度、風險門檻與基準率歷史，必要時把決策降級；它不負責整合四個領域的投資觀點。

\subsection{資料收集}

\subsubsection{四類資料與時間控制}
樣本是九檔美國大型股（AAPL、\allowbreak NVDA、\allowbreak GOOGL、\allowbreak MSFT、\allowbreak AMZN、\allowbreak JPM、\allowbreak MCD、\allowbreak INTC、\allowbreak GE），分析日為 2021 至 2025 年共 20 個季末。每個案例只用分析日以前就已公開的資料（表~\ref{tab:data}）：財報依「申報日」而不是「會計期末」篩選；總經資料使用當時實際已發布的歷史版本；資料快照以內容的「雜湊值」（可想成這份資料的指紋碼，內容一改就會不同）來辨識，四組決策機制讀的是同一份快照。

\begin{table}[H]
\centering\small
\caption{資料來源與處理方式}
\label{tab:data}
\begin{tabularx}{\textwidth}{@{}L{.10} L{.25} Y@{}}
\toprule
類別 & 來源 & 內容與處理 \\
\midrule
技術面 & 調整後日行情 & 近 20 日報酬、均線相對位置、年化波動度；只用分析日以前的價格 \\
基本面 & SEC 財報 & 營收、淨利、營業現金流，並與去年同期比較；依申報日篩選 \\
情緒面 & 金融新聞標題（FNSPID、Alpha Vantage） & 分析日前 90 天的全部標題（既有資料共 83,093 則），以 FinBERT 逐則評分後彙整 \\
總經面 & FRED／ALFRED & 失業率、利率、物價等，只採分析日當時已發布的數值 \\
\bottomrule
\end{tabularx}
\end{table}

九檔股票主要依「資料取不取得到」挑選，不是從 2021 年的成分股中隨機抽出。實驗規則記錄了兩項調整：新聞來源沒有涵蓋的 ASTS 被排除；早期新聞有缺口的 LLY 被換成 INTC。所以樣本是有目的地挑選的，也有存活偏誤的風險（只看到活下來的公司）。主要比較保留全部 180 案；先前排除 37 案的分析，只當成新聞來源與覆蓋範圍的診斷，不用它的結果去訂一個統一的新聞數量門檻。

\subsubsection{資料品質與版本風險}
為了讓快照只反映分析日當時能知道的資訊，本研究做了三件事。第一，財報依申報日期而不是會計期末篩選，並且用相同會計科目、相同期間長度與去年同期比較，避免把不同期間混在一起。第二，年增率的分母用前期的絕對值，這樣當基期是負數、狀況又變差時，年增率會是負的，不會因為正負號而誤導。第三，快照以雜湊值識別，同一個模型、同一個案例的四組決策機制，讀的都是同一份快照與鎖定報告。行情使用 yfinance 的 auto\_adjust=True；調整係數反映的是「下載當下」的拆股與股利處理，並不是各分析日當時保存的版本，因此仍有事後被修訂的風險。

在時間上，新聞依發布時間與時區篩選，財報以申報日判斷「何時才看得到」，總經資料用分析日當時已發布的歷史版本，價格輸入只使用截止時間以前的資料。每個案例都會記錄資料版本、快照雜湊值、證據編號、程式版本、模型版本與稽核紀錄。不過，「各組共用同一份快照」並不能單獨證明資料內容是對的，還必須檢查資料的可用時間與來源。模型版本、新聞來源與行情調整版本，都會影響能不能精確重現結果，所以要和快照、執行紀錄一起保存。

\subsection{研究設計選擇與理由}
\begin{table}[H]\centering\small
\caption{主要設計選擇及其用途}
\begin{tabularx}{\textwidth}{@{}L{.17} L{.19} Y@{}}\toprule
設計 & 設定 & 理由 \\\midrule
辯論者 & 看多與看空 & 用兩個相反的方向，建立最簡單的多空論證結構。\\
辯論輪次 & 三輪 & D 組：第一輪提出論證、第二輪自我反駁、第三輪整合修正。C 組：三輪都維持原立場，並依看得到的內容批評與整合。輪數只是操作上的選擇，不宣稱是最佳的。\\
交換時點 & 第二輪 & 先有可以引用的原始主張，再要求從反方向檢查。\\
交換輪能看到的資訊 & 只有自己的第一輪 & 把反駁的對象限定為自己的主張，避免變成模仿對手。\\
第三輪 & 回到原立場 & 把反面證據與不確定性整合進來，觀察修正後的論點。\\
裁決與驗證 & 裁決者之後接守門員 & 把「決定方向」、「證據是否有效」與「風險控制」分成三件事。\\
主要比較 & C／D 預定呼叫次數相同 & 控制其中一項推論資源的差異；實際用掉的 token 與重試成本另外報告。\\
分析頻率 & 每季一次 & 配合財報按季更新的節奏，並沿用既有資料流程與 20 個決策時間點。\\
主要期間 & 60 個交易日 & 當作季度研究的實際持有期間；30 與 90 日用來檢查結果對期間敏不敏感，不宣稱 60 日在理論上最好。\\
資訊領域 & 技術、基本、情緒、總經 & 沿用既有四類證據，區分「整理資料」與「做決策」兩層的專業化。\\\bottomrule
\end{tabularx}\end{table}
各組共通的輸出格式與一般性的證據驗證設定，屬於實作上的參數，與 C／D、模型組合、任務架構等「研究要比較的因素」分開記錄。D 組專屬的「自我引用驗證」，以及驗證不過時觸發的重試，則是整套 D 流程的一部分，和立場交換、資訊隔離一起構成要比較的對象。

\subsection{實驗設計}

\subsubsection{共同的實驗設定}
各條件使用相同的案例、資料快照、結果（outcome）與回測。研究代理人用的模型與輸出、中立報告與證據登錄都固定，避免更換辯論者時，同時改變了輸入資訊。C、D 兩組的三輪辯論與裁決者呼叫次數對齊；模型設定、提示詞指紋碼、實際的隨機種子、token 用量與重試次數，都記錄在稽核紀錄中。

「記憶」只包含分析日以前已經「到期」（60 日結果已經揭曉）的同一檔股票案例，禁止讀取該案例本身的未來報酬，也禁止讀取其他條件的測試結果。RQ1 目前的設定是：C、D 的裁決者可以讀這類已到期記憶，A、B 不行（見\frref{sec:memory}）。依研究服務程式讀取記憶的規則，系統只取「相同實驗規則雜湊值、同一檔股票、同一組別、相同期間與成本模型」的過去案例，且該案例的到期日與分析日都必須早於目前案例的分析日，只取經守門員處理後的決策、排除降級案例，每組最多取最近 20 筆。也就是說，C 與 D 各自只讀自己這一組過去的決策與對錯，彼此隔離，不會讀到另一組的決策；D 的記憶內容是 D 自己過去決策的結果，和 C 不同，這也是「整套流程」差異的一部分。

須注意的是，以上記憶的時間與條件邊界已於 2026 年 10 月 15 日由封存快照核對，但其語意影響與原執行環境的完整重現仍有限制（見\frref{sec:repro}）；在核對完成前，本文不把 C、D 的差異解讀為單一機制的因果效果。RQ2、RQ3 若執行，必須讓各條件的記憶彼此隔離：每個條件只讀自己已到期的案例，不共享決策，並逐案記錄當時可讀內容與到期規則。由於後面案例的輸入會受到前面決策的影響，比較的是「整套流程動態運作的效果」，而不是各次獨立判斷的差異。模型服務的版本與內容被截斷的情形，也必須記錄。

\subsubsection{四組決策機制}
A 組（單次判斷）與 B 組（七次獨立投票）只當作補充對照，主要比較是 C 與 D。四組的流程如表~\ref{tab:arms}。\grp{B}、\grp{C}、\grp{D} 三組的決策呼叫次數對齊（各 7 次），\grp{A} 則是單次判斷的基準。這裡的 7 次是「預定」的次數：D 組若引用自己原文的驗證沒通過，可能另外觸發重試；呼叫次數相同，不代表實際的重試次數、token 用量或運算量相同，而且辯論組因為要讀歷史論述，輸入也比較長。每次決策都輸出買進（做多）、持有（不進場）或賣出（放空），並附上 60 個交易日的預期報酬、信心度、理由與引用的證據。本文用 A、B、C、D 代表決策機制；RQ2 的基礎模型則另外用 M1、M2 表示，以免搞混。

\begin{table}[H]
\centering\small
\caption{四組決策機制}
\label{tab:arms}
\begin{tabularx}{\textwidth}{@{}c L{.17} c Y@{}}
\toprule
組別 & 名稱 & 呼叫次數 & 流程 \\
\midrule
\grp{A} & 單次判斷 & 1 & 讀完報告後直接決策 \\
\grp{B} & 獨立投票 & 7 & 同樣的提示詞獨立取樣七次，以多數決決定 \\
\grp{C} & 固定立場辯論 & 7 & 看多與看空辯論三輪：第一輪各自提出論證；第二、三輪可以看到已完成輪次中雙方的論述；最後由裁決者決策 \\
\grp{D} & 立場交換辯論 & 7 & 第二輪雙方換邊，只能看到自己的第一輪並且必須反駁它；第三輪換回原立場；最後由裁決者決策 \\
\bottomrule
\end{tabularx}
\end{table}

\subsubsection{固定立場辯論 C}

C 組的看多與看空在三輪中都維持原本的立場，每輪各呼叫一次，依規定可以看到的內容進行批評與整合，最後交給固定的裁決者；預定決策呼叫次數共七次。

\subsubsection{立場交換辯論 D}
\paragraph{比較的對象。} 本專題估計的是「整套 D 流程」相對於「整套 C 流程」的效果差，其中包含立場交換、交換輪看不到對手、強制引用自己原文，以及相對應的驗證與重試。因為這些因素同時改變，所以不能把結果歸因於單獨的「立場交換」。

\paragraph{立場交換辯論的流程。} 表~\ref{tab:rounds} 是 D 組的三輪流程。第二輪刻意不讓角色看到對手的論述：如果同時看得到對手，模型容易變成在回應對手，而不是檢討自己的主張。辯論組的裁決者必須逐項核對雙方引用的證據，並綜合雙方論點，在買進與賣出之間選一個。

\begin{table}[H]
\centering\small
\caption{立場交換辯論的三輪流程}
\label{tab:rounds}
\begin{tabularx}{\textwidth}{@{}L{.11} Y Y@{}}
\toprule
輪次 & 內容 & 目的 \\
\midrule
第 1 輪 & 看多與看空代理人各自提出論點並引用證據 & 建立正反兩方的初始立場 \\
第 2 輪 & 雙方換邊，只能看到自己第一輪的論述，必須引用並反駁自己原本的主張 & 逼代理人找出自己論點的弱點 \\
第 3 輪 & 雙方換回原立場，整合前兩輪的論點與反駁 & 產生整合後的論點，並說明不確定的地方 \\
裁決 & 讀取完整辯論與報告，逐項核對證據後擇一方向 & 避免只看誰講得比較大聲 \\
\bottomrule
\end{tabularx}
\end{table}


在三輪之中，模型本身固定不變：只有「立場」在第二輪反過來、第三輪再換回。「必須引用自己第一輪原文」的驗證規則與重試，是處理流程的一部分，正式實驗前會保存驗證程式的指紋碼。目前的原文核對，是要求至少 70\% 的字詞連續相同；這只能檢查引用有沒有來源，不能證明反駁的內容真的有道理。

\subsubsection{RQ1 的設計}
在相同的案例資料與相同的預定決策呼叫次數下，自我反駁式立場交換辯論（D）的 60 日方向準確率與年化 Sharpe，是否高於固定立場辯論（C）？

RQ1 維持「同一個模型當辯論者」、同樣的資料與七次預定呼叫，估計整套 D 相對於整套 C 的方向準確率差與 Sharpe 差；A、B 只作補充對照。股票代號與公司名稱在輸入模型時都做了匿名化，並且鎖定版本、事先寫好分析計畫，讓執行規則前後一致。匿名化並不能排除模型從數字、日期與產品線索猜出公司。此外，這 180 案先前已用於前導試跑與規則調整，所以這個比較是對既有樣本的事後回頭分析，不能當成對「全新樣本」的獨立驗證。

\subsubsection{RQ2 的模型組合設計}
\frlabel{sec:rq2design}

RQ2 只改變「多空辯論者用哪個模型」，其餘全部固定：案例、資料快照、中立報告、研究代理人、三輪規則、立場交換規則、裁決者、守門員、結果與回測。我們用 M1、M2 代表兩個基礎模型：M1 預定沿用既有主分析的 qwen3:32b；M2 已指定為 gemma3:27b，其確切權重、版本與推理設定（包括裁決者與研究摘要用的模型），會在執行前寫入實驗規則並鎖定。同一模型的條件是 M1/M1；不同模型的條件有兩種平衡配置：配置 X（看多 M1、看空 M2）與配置 Y（看多 M2、看空 M1）。兩種配置都要有 C 與 D，避免 M1 永遠擔任看多方。M2/M2 的同模型條件是必要組別（不再視資源而定），用來分開「模型本身的能力」與「組合的效果」。每個條件至少以三個不同的隨機種子重複執行（確切數量於登記前鎖定），並報告種子之間的變異，避免把單次生成的雜訊當成效果。模型名稱標籤與論點的呈現順序也要控制，以降低文風與裁決偏好帶來的影響。

\begin{table}[H]\centering\small
\caption{不同模型的 D 組中，模型身分與立場的對應（模型不會隨立場一起交換）}
\begin{tabularx}{\textwidth}{@{}lYY@{}}\toprule
輪次 & 配置 X & 配置 Y\\\midrule
第 1 輪 & M1 看多；M2 看空 & M2 看多；M1 看空\\
第 2 輪 & M1 看空；M2 看多 & M2 看空；M1 看多\\
第 3 輪 & M1 看多；M2 看空 & M2 看多；M1 看空\\\bottomrule\end{tabularx}\end{table}

\paragraph{主要比較量。} RQ2 要問的是：「立場交換的效果，會不會隨著模型組合而改變？」所以主要比較量直接定義為「差中之差」：
\[
\Delta_{\mathrm{acc}}=(\mathrm{Acc}_{D,\mathrm{het}}-\mathrm{Acc}_{C,\mathrm{het}})-(\mathrm{Acc}_{D,\mathrm{hom}}-\mathrm{Acc}_{C,\mathrm{hom}}),
\]
（het 代表不同模型的條件，hom 代表同一模型的條件。）Sharpe 的對應比較量為 $\Delta_{\mathrm{SR}}=(\mathrm{SR}_{D,\mathrm{het}}-\mathrm{SR}_{C,\mathrm{het}})-(\mathrm{SR}_{D,\mathrm{hom}}-\mathrm{SR}_{C,\mathrm{hom}})$。兩者衡量的都是立場交換的效果是否隨模型組合而變。同模型的對照有 M1/M1 與 M2/M2 兩個，上式分別以兩者當作 hom，各得一組結果；若兩組結果方向不一致，結論只限於各自的模型組合。M1/M1 與 M1/M2 的比較，同時也改變了模型的能力組成，所以結果只適用於這組指定的模型，不能解讀成「異質模型一般都有好處」。

\paragraph{兩種配置如何合併（預定做法，執行前鎖定）。} 在不同模型的條件下，每個案例的方向正確性，取配置 X 與 Y 的平均（0、0.5 或 1）；日報酬序列則取兩種配置組合日報酬的等權平均，再用與 RQ1 相同的規則算 Sharpe。兩種配置共用同樣的案例與結果，不能當成兩倍的獨立樣本；也不能把兩個 Sharpe 直接平均，或把兩條序列接在一起。統計推論以「分析日」為單位，把同一天的案例綁在一起做重抽樣（區塊 bootstrap），估計上述兩個比較量的信賴區間與 $p$ 值。目前分析原型把 M1/M1 與 M2/M2 兩個基準各自的準確率、Sharpe 共四個端點放在同一個 Holm 校正 family，採雙尾檢定；這是正式登記前須確認並鎖定的提案，不從兩個基準中挑選較好的結果。由於平均後的正確性不再是「對／錯」二元結果，所以不使用一般的 McNemar 檢定。

\subsubsection{RQ3 的領域分工設計}
\frlabel{sec:rq3design}
通才（Generalist）架構是由共同的看多／看空代理人，讀四類資訊整合後的研究報告。領域專業化架構則把同一份研究輸出依領域切開，分給技術面、基本面、情緒面與總經面各自的多空辯論；每個領域做出一份「領域看法」（Domain View），再由最終整合者（Final Aggregator）整合成一個候選決策，之後接相同的守門員與回測。RQ3 比較的是「整套領域化架構」相對於整套通才架構；其中同時包含決策層的任務拆分、資訊切分與額外的呼叫次數，所以結論只稱為架構效果，不稱為單純的分工效果。預定的流程如表~\ref{tab:rq3flow}。為固定辯論規則、避免再混入立場交換，兩種架構都採用 RQ1 的固定立場規則 C；呼叫數與資訊切分仍同時改變，因此 RQ3 比較的是架構套組。

\begin{table}[H]\centering\small
\caption{RQ3 預定流程（設計階段，尚未執行）}
\label{tab:rq3flow}
\begin{tabularx}{\textwidth}{@{}L{.17} Y Y@{}}\toprule
項目 & 通才組 & 領域專業化組 \\\midrule
辯論者與證據 & 共同的看多／看空，讀整合報告 & 每個領域各有看多／看空，只讀該領域的證據 \\
辯論規則 & C：三輪固定立場 & 每個領域都用 C：三輪固定立場 \\
領域內裁決 & 無 & 每個領域有一位裁決者，輸出標準格式的領域看法（方向、支持證據、反面證據、不確定性、資訊缺口；資訊不足時可以棄權） \\
上層整合 & 裁決者直接裁決 & 最終整合者讀四份領域看法，用「裁決」（不是投票）決定單一候選方向 \\
後處理 & 相同的守門員與回測 & 相同的守門員與回測 \\
預定決策呼叫次數（示例） & 7 & $4\times7+1=29$ \\
\bottomrule\end{tabularx}\end{table}

表中的呼叫次數只是這個設計下的預定示例，不是實際測到的用量。真正做研究時，研究、辯論、整合與重試的呼叫次數、輸入／輸出 token、推理時間與整體延遲都要記錄，用來呈現「績效」與「成本」之間的取捨。四份領域看法都針對同一個 60 交易日的方向；所謂 4:0、3:1、2:2 的意見分歧，只適用於四個都有效的二元方向，資訊不足或棄權的另外列出。這裡的「領域專業化」，指的是角色與證據範圍上的分工，並不是已經證明模型具有金融專業能力，也不是對模型做金融微調。

RQ3 主要比較的是整套架構的績效與成本，結論只限於「架構效果」：通才組與領域專業化組同時改變了分工、呼叫次數與資訊切分，績效差距不能說成「純粹是分工造成」。為了分辨架構效果的來源，另外設兩個必要的對照：（一）運算資源對齊對照，讓通才組使用與領域化架構相同的總預定呼叫次數（約 29 次，例如多次取樣後投票）；（二）資訊切分消融，把資訊同樣切成四個領域，但各領域不辯論，改由單一裁決者直接讀四份領域證據。目前原型將運算對照設為 28 次彼此不可見的獨立提案與 1 次整合裁決；切分對照為單一裁決者直接讀取四份領域證據、不進行辯論。兩者均使用相同的守門員與回測，並有逐步續跑的合成測試；這些設定仍須在真實模型資格檢驗後寫入正式登記。

主要指標是方向準確率與年化 Sharpe，怎樣算勝出、多個比較如何校正，都在執行前事先寫定（表~\ref{tab:rqcompare}）。

此外，目前另有隔離 pilot／研究輸入庫執行入口，以及從唯讀凍結快照建立完整四組共同案例面板的離線分析入口。RQ3 分析原型對領域組分別減去 C、29 次通才對照與資訊切分對照，產生三個比較的準確率差與 Sharpe 差，提供六端點共同 Holm 校正的保守提案；若正式計畫改用主要／次要 family，必須在執行前另行鎖定並驗證。分析程式已用合成的完整 180 案快照測試，尚無真實 RQ3 結果，也不能以此取代模型資格。RQ2 的研究輸入庫另用不在正式種子清單中的 seed（例如 904），在任何決策前鎖定並暫停；正式 905／906／907 共用該逐案報告，避免來源工作佔用正式案例。

\subsubsection{評估方式}

\paragraph{回測與主要比較。} 在分析日後第一個交易日的開盤價進場，持有 60 個交易日後以開盤價出場，並扣除用 Corwin--Schultz 方法估計的買賣價差成本 \citep{corwin2012spread}。主要比較是 \grp{D} 對 \grp{C}，用兩個共同的主要指標：方向準確率（用配對的 McNemar 檢定 \citep{mcnemar1947}，這是檢查兩種方法對同一批案例的對錯有沒有系統性差異的檢定），以及年化 Sharpe 比率的差（用 Ledoit--Wolf 檢定 \citep{ledoit2008sharpe}）；兩個指標一起以 Holm 法 \citep{holm1979} 校正，採雙尾檢定。只有當兩個指標的差異都是正的，而且各自校正後 $p<0.05$，才判定 \grp{D} 比較好。這是「兩個指標都要顯著為正」的交集規則，搭配雙尾檢定，只承認 \grp{D} 的正向結果；在這種規則下，Holm 校正偏保守。保留這套規則，是因為它在看到結果之前就已鎖定，不因結果而更動。RQ1 的問句寫「是否高於」，但檢定採雙尾，是為了同時報告負向的差異。另外，30 與 90 日、不計成本、排除先前新聞來源／覆蓋診斷的子集（37 案；名單與選取規則須在定稿前核對並附上）、以及依日期分組的信賴區間，都屬於敏感性分析（換個設定看結論會不會變）。

\paragraph{報酬與 Sharpe 的算法。} 日報酬先依每個案例「進場開盤價、每日收盤價、出場開盤價」算出部位的價值變化，並扣掉進出場的價差成本；同一天的各個案例用固定權重 1/180，沒有持有或不在交易期間的報酬記為零。年化 Sharpe 為 $\sqrt{252}\bar r/s(r)$，也就是「平均日報酬除以日報酬的標準差，再乘上一年的交易日數開根號」。其中，案例數 $n_{\mathrm{case}}=180$，分析日的分組數為 20，$T$ 是對齊後日資料的交易日筆數（2026 年 10 月 15 日由唯讀快照重算：2021-04-01 至 2026-03-31，$T=1219$），這三個數字不能互相替換；標準差的分母用 $T-1$，無風險利率明設為零。累積報酬為 $\prod_t(1+r_t)-1$，不是各案例 60 日報酬的平均；最大回撤（資產從最高點下跌的最大幅度）也由這條日序列計算，這種低曝險配置只適合用來比較各組，不適合當成實際績效。方向是否正確，是依未扣成本的開盤價漲跌判斷，和扣成本後淨報酬的正負不同。Ledoit--Wolf 檢定的重抽樣區塊長度為 20 日，每次重抽樣的長度為 $20\lfloor T/20\rfloor$，共重抽 1,999 次，隨機種子 905；qwen 的基礎推理種子是 905、溫度 0.2（投票組 0.8），各次呼叫實際用的種子另外保存在稽核紀錄中。

\paragraph{次要指標。} 次要指標包含累積報酬、最大回撤、決策分布、信心度、扣除交易成本後的報酬，以及意見分歧。RQ2 另外報告看多與看空雙方的分歧、不同模型決策的差異，以及各模型偏向看多或看空的傾向；RQ3 報告四個領域各自的立場、一致程度，以及跨領域衝突如何處理。要注意：意見分歧增加、引用更完整，或比較常看空，都不直接代表績效有進步；信心度是模型自己報的數字，不能自動當成校準過的機率。候選決策層與守門員之後的結果分開報告覆蓋率、有出手時的準確率與日報酬；無風險利率設為零與低曝險配置的限制，也要揭露。

\paragraph{機制指標（預先定義）。} 為了檢查「自我反駁」是否真的改變了推理，而不只是改變文字，定義三個機制指標：（一）反面證據引用率，即第二輪引用的證據中，原本支持對立立場者所占的比例；（二）裁決者引用證據的組成，即裁決者引用的證據中，支持買進與支持賣出各占的比例；（三）方向翻轉率，即第三輪預期報酬的方向與第一輪不同的比例。RQ2、RQ3 會在執行前把這些指標寫入登記；RQ1 的 180 案先前沒有登記這些指標，只能在原始輸出補齊後以探索性方式描述，不作為確認性結論。這些指標只有在同時伴隨績效改善時，才被解讀為機制的證據。

\subsubsection{統計分析}

RQ1 沿用配對的 McNemar 檢定、Ledoit--Wolf 檢定與區塊重抽樣、雙尾檢定、兩個指標共同的 Holm 校正，以及依日期分組的敏感性分析；只有當兩個指標的差異都是正的、而且校正後仍顯著，才算支持 D 比較好。180 個案例共用 20 個分析日、同一檔股票反覆出現，持有期間也有重疊，所以「各案例互相獨立」這個假設必須用分組／區塊的方法檢查；而且日期分組的數目有限，不能把敏感性分析當成已經充分校正的證明。由於 McNemar 檢定把 180 案視為互相獨立，本文在解讀 RQ1 時，把依日期分組的信賴區間與預定的檢定並列，作為考慮案例相依之後的判斷依據。兩個指標的區間都包含 0（準確率差 $[-7.2,+0.6]$ 個百分點、Sharpe 差 $[-0.47,+0.36]$），與預定檢定的結論一致，所以結論不隨檢定方式而改變。這是看到結果之後才提高日期分組推論的地位，並沒有更改預定的判準；RQ2、RQ3 則從一開始就以依日期分組的推論為主分析。結果沒有顯著，不等於兩者一樣好或「不比較差」，而「算作一樣好」的界線也不能看完結果才補訂。

RQ2 的主要檢定是\frref{sec:rq2design}的差中之差 $\Delta_{\mathrm{acc}}$ 與 $\Delta_{\mathrm{SR}}$，而不是分別檢定各條件下的 D$-$C。如果一組顯著、另一組不顯著，這件事本身並不能證明「效果隨模型組合而變」。兩種配置共用案例，平均後的正確性不再是二元結果，所以不能直接套用一般 McNemar；Sharpe 也要由配對的日序列計算，不能把兩個 Sharpe 平均，或把兩種配置接成兩倍的觀測數。

RQ3 在相同的「股票—日期」案例上，比較架構與成本；各領域的論點只是同一案例的中間輸出，不是四筆獨立的金融樣本。主要指標與探索性的診斷要分開呈現，統計推論也要處理共同日期、同一檔股票與重疊持有期造成的相依，避免只挑對自己有利的次要指標。

三個研究問題比較的東西不同，不能把 RQ1 的判準不加說明地套到 RQ2、RQ3，整理於表~\ref{tab:rqcompare}；其中 RQ2、RQ3 兩欄是預定的設計，執行前會鎖定。共同的回測規則（進出場、成本、權重、無風險利率）見上面的「評估方式」。

\begin{table}[H]\centering\small
\caption{三個研究問題的比較設計}
\label{tab:rqcompare}
\begin{tabularx}{\textwidth}{@{}L{.13} Y Y Y@{}}\toprule
 & RQ1 & RQ2 & RQ3 \\\midrule
比較條件 & D 對 C（同一模型 qwen3:32b） & 同模型 M1/M1 與不同模型 M1/M2（平衡配置）下的 D 對 C & 領域專業化對通才 \\
主要比較量 & 方向準確率差；Sharpe 差 & 差中之差 $\Delta_{\mathrm{acc}}$、$\Delta_{\mathrm{SR}}$ & 方向準確率差；Sharpe 差；成本另外報告 \\
分析單位 & 180 個案例（Sharpe 由日序列計算） & 同左；不同模型的條件先合併 X、Y 兩配置 & 同左 \\
統計方法 & 配對 McNemar；Ledoit--Wolf 檢定與區塊重抽樣 & 以日期為單位分組的區塊重抽樣 & 同 RQ2，細節執行前鎖定 \\
判準 & 兩個指標都為正，且 Holm 校正後 $p<0.05$ & 兩個基準各兩指標，共四端點 Holm 提案；雙尾 & 三比較共六端點 Holm 提案；主比較兩指標都須為正且校正顯著 \\
成本 & 預定呼叫次數對齊；token 與重試另外報告 & 預定呼叫次數對齊 & 呼叫次數、token、延遲全部報告，不併入門檻 \\
狀態 & 已完成（回溯比較） & 設計階段 & 設計階段 \\
\bottomrule\end{tabularx}\end{table}

\subsection{工具及設備}

\subsubsection{既有工具與 GPU}
\begin{table}[H]
\centering\small
\caption{工具及設備}
\label{tab:tools}
\begin{tabularx}{\textwidth}{@{}L{.17} Y Y@{}}
\toprule
項目 & 工具 & 用途 \\
\midrule
程式語言與框架 & Python、FastAPI & 後端服務與網頁介面 \\
流程控制 & LangGraph & 管理研究、決策、風控與回測的狀態流程 \\
模型呼叫 & Ollama、LiteLLM & 部署開源模型 qwen3:32b，並提供統一的模型呼叫介面 \\
資料來源 & yfinance、\allowbreak SEC EDGAR、\allowbreak FRED／ALFRED、\allowbreak FNSPID、\allowbreak Alpha Vantage & 行情、財報、總經與新聞 \\
情緒分析 & FinBERT & 新聞標題的正負面評分 \\
儲存與部署 & SQLite、Docker & 保存案例與實驗紀錄；容器化部署 \\
運算資源 & 雲端 GPU 主機 & 執行模型推理 \\
\bottomrule
\end{tabularx}
\end{table}
